\chapter{Réévaluation à base de règles}
\label{cha:reevaluation-base-regles}

Dans le chapitre~\ref{cha:etat-art-modelisation-parametrique}, nous avons passé
en revue différents systèmes de modélisation, leurs méthodes pour la
réévaluation de spécifications paramétriques et les différentes solutions qu'ils
proposent pour la nomination persistante des entités topologiques. %
Nous avons constaté que la plupart des systèmes ne sont adaptés qu'à la
conception et la réévaluation de modèles paramétriques d'une dimension donnée
(généralement en 3D) et définissent des mécanismes de nommage hétérogènes en
fonction de la dimension des cellules. %
De ce fait, leurs méthodes ne sont définies que pour un certain ensemble
d'entités topologiques. %
Ces méthodes sont ainsi difficiles à étendre vers d'autres dimensions et donc
non généralisables. %

Dans ce chapitre, nous nous appuyons sur les travaux de Cardot~\cite{cardot2019}
et les concepts qu'elle propose pour mettre en \oe uvre un système de
réévaluation générique grâce à l'utilisation des cartes généralisées et des
règles Jerboa. %
En particulier, nous nous appuyons sur ses contributions pour l'implantation d'un système de
nomination persistante en enrichissant la reconstitution des historiques d'évolution des
orbites à l'aide du suivi des orbites et de leurs origines. %
Cette amélioration participe à l'amélioration de la procédure
d'appariement des entités topologiques ainsi qu'à la définition de stratégies
accessibles aux utilisateurs pour la réévaluation. %

Nous commençons avec une première section dans laquelle nous
présentons deux cas d'études de réévaluation. %
Dans une deuxième section, nous présentons notre approche de la
nomination persistante avec, d'une part, la nomination des brins grâce aux
règles Jerboa et, d'autre part, la nomination des orbites grâce à la
formalisation des évènements (voir
chapitre~\ref{cha:formalisation-evenements}). %
Dans une troisième section, nous présentons notre mécanisme d'appariement des
orbites lors du processus de réévaluation. %
Nous proposons, ensuite, un système de stratégie pour paramétrer la réévaluation
d'une spécification paramétrique. %
Enfin, dans la quatrième section, nous donnons un état des lieux de ce qui a été
implanté avant de conclure le chapitre. %


\section{Cas d'études}
\label{sec:reevaluation-base-regles-cas-etude}


Tout au long de ce chapitre, nous présentons nos différentes approches pour la
réévaluation, notamment la nomination persistante des entités topologiques et
leurs appariements. Nous nous appuyons, pour cela, sur deux exemples de constructions et leurs réévaluations, qui
nous servent de cas d'études. %
Pour chacune de ces constructions, nous proposons les représentations
géométriques (par étape de construction) accompagnées de leurs spécifications
paramétriques. %

\subsection{Construction nº~1}
\label{sec:reevaluation-base-regles-premier-cas}

\begin{sidewaysfigure}
	\centering
	\begin{subfigure}[t]{0.1\textwidth}
		\includegraphics[width=\textwidth]{1-evaluation-square}
		\caption{}
		\label{fig:reevaluation-base-regles-cs1-1-evaluation-square}
	\end{subfigure}
	\hfill
	\begin{subfigure}[t]{0.1\textwidth}
		\includegraphics[width=\textwidth]{2-evaluation-extrusion}
		\caption{}
		\label{fig:reevaluation-base-regles-cs1-2-evaluation-extrusion}
	\end{subfigure}
	\hfill
	\begin{subfigure}[t]{0.1\textwidth}
		\includegraphics[width=\textwidth]{3-evaluation-insert}
		\caption{}
		\label{fig:reevaluation-base-regles-cs1-3-evaluation-insert}
	\end{subfigure}
	\hfill
	\begin{subfigure}[t]{0.1\textwidth}
		\includegraphics[width=\textwidth]{4-evaluation-extrusion}
		\caption{}
		\label{fig:reevaluation-base-regles-cs1-4-evaluation-extrusion}
	\end{subfigure}
	\hfill
	\begin{subfigure}[t]{0.1\textwidth}
		\includegraphics[width=\textwidth]{5-evaluation-triangulation}
		\caption{}
		\label{fig:reevaluation-base-regles-cs1-5-triangulation}
	\end{subfigure}
	\hfill
	\begin{subfigure}[t]{0.1\textwidth}
		\includegraphics[width=\textwidth]{6-evaluation-collapse}
		\caption{}
		\label{fig:reevaluation-base-regles-cs1-6-evaluation-collapse}
	\end{subfigure}
	\hfill
	\begin{subfigure}[t]{0.1\textwidth}
		\includegraphics[width=\textwidth]{7-evaluation-chamfer}
		\caption{}
		\label{fig:reevaluation-base-regles-cs1-7-evaluation-chamfer}
	\end{subfigure}
	\hfill
	\begin{subfigure}[t]{0.1\textwidth}
		\includegraphics[width=\textwidth]{8-evaluation-colour}
		\caption{}
		\label{fig:reevaluation-base-regles-cs1-8-evaluation-color}
	\end{subfigure}
	\hfill
	\begin{subfigure}[t]{0.1\textwidth}
		\includegraphics[width=\textwidth]{9-evaluation-color}
		\caption{}
		\label{fig:reevaluation-base-regles-cs1-9-evaluation-color}
	\end{subfigure}
	\caption{Évaluation cas nº~1~:
		(\subref{fig:reevaluation-base-regles-cs1-1-evaluation-square})
		\texttt{1-\allowbreak{}creation\-Carre(pos)}~;
		(\subref{fig:reevaluation-base-regles-cs1-2-evaluation-extrusion})
		\texttt{2-\allowbreak{}extrusion\-Face(PN\textsubscript{6},vec)}~;
		(\subref{fig:reevaluation-base-regles-cs1-3-evaluation-insert})
		\texttt{3-\allowbreak{}insertion\-Sommet(PN\textsubscript{46})}~;
		(\subref{fig:reevaluation-base-regles-cs1-4-evaluation-extrusion})
		\texttt{4-\allowbreak{}extrusion\-Face\-Volume(PN\textsubscript{47},vec)}~;
		(\subref{fig:reevaluation-base-regles-cs1-5-triangulation})
		\texttt{5-\allowbreak{}tri\-an\-gu\-la\-t\-i\-on\-Face(PN\textsubscript{33})}~;
		(\subref{fig:reevaluation-base-regles-cs1-6-evaluation-collapse})
		\texttt{6-\allowbreak{}contraction\-Face(PN\textsubscript{63})}~;
		(\subref{fig:reevaluation-base-regles-cs1-7-evaluation-chamfer})
		\texttt{7-\allowbreak{}chanfrein\-Sommet(PN\textsubscript{55})}~;
		(\subref{fig:reevaluation-base-regles-cs1-8-evaluation-color})
		\texttt{8-\allowbreak{}co\-lo\-ra\-t\-i\-on\-Face(PN\textsubscript{35},Bleu)}~;
		(\subref{fig:reevaluation-base-regles-cs1-9-evaluation-color})
		\texttt{9-\allowbreak{}co\-lo\-ra\-t\-i\-on\-Face(PN\textsubscript{68},
			Rouge)}}
	\label{fig:reevaluation-base-regles-cs1-evaluation}

	\vspace{1cm}

	\begin{subfigure}[t]{0.1\textwidth}
		\includegraphics[width=\textwidth]{1-reevaluation-square}
		\caption{}
		\label{fig:reevaluation-base-regles-cs1-1-reevaluation-square}
	\end{subfigure}
	\begin{subfigure}[t]{0.1\textwidth}
		\includegraphics[width=\textwidth]{1ADD-reevaluation-insert}
		\caption{}
		\label{fig:reevaluation-base-regles-cs1-add-10-reevaluation-insert}
	\end{subfigure}
	\hfill
	\begin{subfigure}[t]{0.1\textwidth}
		\includegraphics[width=\textwidth]{2-reevaluation-extrusion}
		\caption{}
		\label{fig:reevaluation-base-regles-cs1-2-reevaluation-extrusion}
	\end{subfigure}
	\hfill
	\begin{subfigure}[t]{0.1\textwidth}
		\includegraphics[width=\textwidth]{4-reevaluation-extrusion}
		\stepcounter{subfigure}\caption{}
		\label{reevaluation-base-regles-cs1-4-reevaluation-extrusion}
	\end{subfigure}
	\hfill
	\begin{subfigure}[t]{0.1\textwidth}
		\includegraphics[width=\textwidth]{5-reevaluation-triangulation}
		\caption{}
		\label{fig:reevaluation-base-regles-cs1-5-reevaluation-triangulation}
	\end{subfigure}
	\hfill
	\begin{subfigure}[t]{0.1\textwidth}
		\includegraphics[width=\textwidth]{6-reevaluation-collapse}
		\caption{}
		\label{fig:reevaluation-base-regles-cs1-6-reevaluation-collapse}
	\end{subfigure}
	\hfill
	\begin{subfigure}[t]{0.1\textwidth}
		\includegraphics[width=\textwidth]{7-reevaluation-chamfer}
		\caption{}
		\label{fig:reevaluation-base-regles-cs1-7-reevaluation-chamfer}
	\end{subfigure}
	\hfill
	\begin{subfigure}[t]{0.1\textwidth}
		\includegraphics[width=\textwidth]{8-reevaluation-colour}
		\caption{}
		\label{fig:reevaluation-base-regles-cs1-8-reevaluation-color}
	\end{subfigure}
	\hfill
	\begin{subfigure}[t]{0.1\textwidth}
		\includegraphics[width=\textwidth]{9-reevaluation-color}
		\caption{}
		\label{fig:reevaluation-base-regles-cs1-9-reevaluation-color}
	\end{subfigure}
	\caption{Réevaluation cas nº~1~:
		(\subref{fig:reevaluation-base-regles-cs1-1-reevaluation-square})
		\texttt{1-carre(pos)}~;
		(\subref{fig:reevaluation-base-regles-cs1-add-10-reevaluation-insert})
		\texttt{ADD~1-\allowbreak{}insertion\-Sommet(3)}~;
		(\subref{fig:reevaluation-base-regles-cs1-2-reevaluation-extrusion})
		\texttt{2-\allowbreak{}extrusion\-Face(PN\textsubscript{6},vec)}~;
		% (\subref{})
		\texttt{DELETE~3-\allowbreak{}insertion\-Sommet(PN\textsubscript{46})}~;
		(\subref{reevaluation-base-regles-cs1-4-reevaluation-extrusion})
		\texttt{4-\allowbreak{}extrusion\-Face\-Volume(PN\textsubscript{47},vec)}~;
		(\subref{fig:reevaluation-base-regles-cs1-5-reevaluation-triangulation})
		\texttt{5-\allowbreak{}tri\-an\-gu\-la\-t\-i\-on\-Face(PN\textsubscript{33})}~;
		(\subref{fig:reevaluation-base-regles-cs1-6-reevaluation-collapse})
		\texttt{6-\allowbreak{}contraction\-Face(PN\textsubscript{63})}~;
		(\subref{fig:reevaluation-base-regles-cs1-7-reevaluation-chamfer})
		\texttt{7-\allowbreak{}chanfrein\-Sommet(PN\textsubscript{55})}~;
		(\subref{fig:reevaluation-base-regles-cs1-8-reevaluation-color})
		\texttt{8-\allowbreak{}co\-lo\-ra\-t\-i\-on\-Face(PN\textsubscript{35},
			Bleu)}~;
		(\subref{fig:reevaluation-base-regles-cs1-9-reevaluation-color})
		\texttt{9-\allowbreak{}co\-lo\-ra\-t\-i\-on\-Face(PN\textsubscript{68},
			Rouge)}}
	\label{fig:reevaluation-base-regles-cs1-reevaluation}
\end{sidewaysfigure}

La première construction, illustrée
figure~\ref{fig:reevaluation-base-regles-cs1-evaluation}, sera notre principal
fil rouge pour ce chapitre. %
Celui-ci va nous permettre d'illustrer nos méthodes de nomination persistante
et de réévaluation. %
Nous détaillons, ci-après, sa spécification paramétrique contenue dans le
listing~\ref{lst:reevaluation-regles-specification-parametrique-init-cas-1}. %

L'application de l'opération, \texttt{1-\allowbreak{}creation\-Carre}, crée une face carrée
(figure~\ref{fig:reevaluation-base-regles-cs1-1-evaluation-square}). %
L'opération, \texttt{2-\allowbreak{}extrusion\-Face}, est appliquée sur cette face,
incidente au brin \(6\), et l'extrude en un cube
(figure~\ref{fig:reevaluation-base-regles-cs1-2-evaluation-extrusion}). %
L'opération, \texttt{3-\allowbreak{}insertion\-Sommet}, est appliquée sur l'arête incidente
au brin \(46\) et y insère un sommet
(figure~\ref{fig:reevaluation-base-regles-cs1-3-evaluation-insert}). %
L'opération, \texttt{4-\allowbreak{}extrusion\-Face\-Volume}, est appliquée sur la face
du dessus, incidente au brin \(47\), et l'extrude en un autre cube
(figure~\ref{fig:reevaluation-base-regles-cs1-4-evaluation-extrusion}). %
L'opération, \texttt{5-\allowbreak{}triangule\-Face}, est appliquée sur la face avant,
incidente au brin \(33\), et la triangule
(figure~\ref{fig:reevaluation-base-regles-cs1-5-triangulation}). %
L'opération, \texttt{6-\allowbreak{}contraction\-Face}, est appliquée sur la face
incidente au brin \(63\), et la contracte en un sommet
(figure~\ref{fig:reevaluation-base-regles-cs1-6-evaluation-collapse}). %
L'opération, \texttt{7-\allowbreak{}chanfrein\-Sommet}, est appliquée sur ce même sommet,
incident au brin \(55\), et fait un chanfrein dessus
(figure~\ref{fig:reevaluation-base-regles-cs1-7-evaluation-chamfer}). %
L'opération, \texttt{8-\allowbreak{}coloration\-Face}, colore la face incidente au brin
\(35\) (figure~\ref{fig:reevaluation-base-regles-cs1-8-evaluation-color}). %
Enfin, la dernière opération, \texttt{9-\allowbreak{}coloration\-Face}, colore la face incidente
au brin \(68\)
(figure~\ref{fig:reevaluation-base-regles-cs1-9-evaluation-color}). %

\begin{lstlisting}[escapechar=\%,
style=script-style,
caption=Spécification paramétrique initiale du cas nº~1,
label=lst:reevaluation-regles-specification-parametrique-init-cas-1,
captionpos=b]
1-%\allowbreak{}creation\-Carre(pos)%
2-%\allowbreak{}extrusion\-Face(\pndart{6}{1n6}, vec)%
3-%\allowbreak{}insertion\-Sommet(\pndart{46}{1n5;2n5})%
4-%\allowbreak{}extrusion\-Face\-Volume(\pndart{47}{1n5;2n6;3n0}, vec)%
5-%\allowbreak{}tri\-an\-gu\-la\-t\-i\-on\-Face(\pndart{33}{1n3;2n2})%
6-%\allowbreak{}contraction\-Face(\pndart{63}{1n4;2n6;4n6})%
7-%\allowbreak{}chanfrein\-Sommet(\pndart{55}{1n5;2n6;3n0;4n4;6n2})%
8-%\allowbreak{}coloration\-Face(\pndart{35}{1n3;2n4;5n0}, Bleu)%
9-%\allowbreak{}coloration\-Face(\pndart{68}{1n5;2n6;3n0;4n4;6n2;7n1}, Rouge)%
\end{lstlisting}

% \xavier{Dans la spéc. paramétrique, parfois tu précises sur quelle entité topologique l'opération a lieu : par exemple, "extrusionFace", parfois non, par exemple, "insertionSommet" $\rightarrow$ sois homogène et précise l'entité pour toutes les opérations.
% \begin{itemize}
% \item carre $\rightarrow$ creationCarre
% \item extrusionFace  $\rightarrow$ OK
% \item extrusionFace\-Volume $\rightarrow$ OK
% \item triangulation $\rightarrow$ triangulationFace (attention, il faut aussi remettre le texte explicatif à jour. En particulier, tu nommes actuellement cette opération "triangulate")
% \item contractionFace  $\rightarrow$ OK
% \item chanfrein $\rightarrow$ chanfreinSommet
% \item coloration $\rightarrow$ colorationFace ou colorationBleuFace
% \item coloration $\rightarrow$ colorationFace ou colorationRougeFace
% \end{itemize}
% }


% \xavier{Même remarque que précédemment pour le renommage des opérations, dans la spéc. param. et dans le texte explicatif.}

Nous procédons ensuite à l'édition de cette spécification paramétrique, dont la
représentation géométrique, à chaque étape, est illustrée dans la
figure~\ref{fig:reevaluation-base-regles-cs1-reevaluation}. %
Cette édition consiste à ajouter une opération d'insertion de sommet
\texttt{9-\allowbreak{}insertion\-Sommet(3)}
(figure~\ref{fig:reevaluation-base-regles-cs1-add-10-reevaluation-insert}) et
à supprimer l'opération \texttt{3-\allowbreak{}insertion\-Sommet} afin que celle-ci ne soit pas
réévaluée. %

% \xavier{Je ne vois pas l'intérêt de la phrase suivante, c'est une redite exacte de la précédente.}
% Les aspects remarquables de la construction et de sa réévaluation que nous
% venons de présenter sont le fait que nous ajoutons une opération et en
% supprimons une autre lors de l'édition de la spécification paramétrique. %


% Nous devrions donc, en toute logique, nous attendre à des changements
% significatifs dans la topologie du modèle réévalué. %



\begin{lstlisting}[escapechar=\%,
style=script-style,
caption=Spécification paramétrique éditée du cas nº~1,
label=lst:reevaluation-regles-specification-parametrique-edit-cas-1,
captionpos=b]
1-%carre(pos)%
ADD 9-%\allowbreak{}insertion\-Sommet(3)%
2-%\allowbreak{}extrusion\-Face(\pndart{6}{1n6}, vec)%
DELETE 3%-\allowbreak{}insertion\-Sommet(\pndart{46}{1n5;2n4})%
4-%\allowbreak{}extrusion\-Face\-Volume(\pndart{47}{1n5;2n6;3n0}, vec)%
5-%\allowbreak{}tri\-an\-gu\-la\-t\-i\-on\-Face(\pndart{33}{1n3;2n2})%
6-%\allowbreak{}contraction\-Face(\pndart{63}{1n4;2n6;4n6})%
7-%chanfrein(\pndart{55}{1n5;2n6;3n0;4n4;6n2})%
8-%coloration(\pndart{35}{1n3;2n4;5n0}, Bleu)%
9-%coloration(\pndart{68}{1n5;2n6;3n0;4n4;6n2;7n1},Rouge)%
\end{lstlisting}


\subsection{Construction nº~2}
\label{sec:reevaluation-base-regles-deuxieme-cas}


\begin{sidewaysfigure}
	\centering
	\begin{subfigure}[t]{0.1\textwidth}
		\centering
		\includegraphics[width=\textwidth]{cs2-1-evaluation-square}
		\caption{}
		\label{fig:reevaluation-base-regles-cs2-1-evaluation-square}
	\end{subfigure}
	\hfill
	\begin{subfigure}[t]{0.1\textwidth}
		\centering
		\includegraphics[width=\textwidth]{cs2-2-evaluation-faceExtrusion}
		\caption{}
		\label{fig:reevaluation-base-regles-cs2-2-evaluatino-extrusion}
	\end{subfigure}
	\hfill
	\begin{subfigure}[t]{0.1\textwidth}
		\centering
		\includegraphics[width=\textwidth]{cs2-3-evaluation-faceCollapse}
		\caption{}
		\label{fig:reevaluation-base-regles-cs2-3-evaluation-collapse}
	\end{subfigure}
	\hfill
	\begin{subfigure}[t]{0.1\textwidth}
		\centering
		\includegraphics[width=\textwidth]{cs2-4-evaluation-chamfer}
		\caption{}
		\label{fig:reevaluation-base-regles-cs2-4-evaluation-chamfer}
	\end{subfigure}
	\hfill
	\begin{subfigure}[t]{0.1\textwidth}
		\centering
		\includegraphics[width=\textwidth]{cs2-5-evaluation-triangulation}
		\caption{}
		\label{fig:reevaluation-base-regles-cs2-5-evaluation-triangulation}
	\end{subfigure}
	\hfill
	\begin{subfigure}[t]{0.1\textwidth}
		\centering
		\includegraphics[width=\textwidth]{cs2-6-evaluation-colourFaceRed}
		\caption{}
		\label{fig:reevaluation-base-regles-cs2-6-evaluation-color}
	\end{subfigure}
	\hfill
	\begin{subfigure}[t]{0.1\textwidth}
		\centering
		\includegraphics[width=\textwidth]{cs2-7-evaluation-colourFaceBlue}
		\caption{}
		\label{fig:reevaluation-base-regles-cs2-7-evaluation-color}
	\end{subfigure}
	\caption{Évaluation cas nº~2~:
		(\subref{fig:reevaluation-base-regles-cs2-1-evaluation-square})
		\texttt{1-carre(pos)}~;
		(\subref{fig:reevaluation-base-regles-cs2-2-evaluatino-extrusion})
		\texttt{2-extrusion(PN\textsubscript{6},vec)}~;
		(\subref{fig:reevaluation-base-regles-cs2-3-evaluation-collapse})
		\texttt{3-\allowbreak{}contraction\-Face(PN\textsubscript{47})}~;
		(\subref{fig:reevaluation-base-regles-cs2-4-evaluation-chamfer})
		\texttt{4-chanfrein(PN\textsubscript{40})}~;
		(\subref{fig:reevaluation-base-regles-cs2-5-evaluation-triangulation})
		\texttt{5-\allowbreak{}tri\-an\-gu\-la\-t\-i\-on\-Face(PN\textsubscript{33})}~;
		(\subref{fig:reevaluation-base-regles-cs2-6-evaluation-color})
		\texttt{6-co\-lo\-ra\-t\-i\-on(PN\textsubscript{12},Rouge)}~;
		(\subref{fig:reevaluation-base-regles-cs2-7-evaluation-color})
		\texttt{7-co\-lo\-ra\-t\-i\-on(PN\textsubscript{57},Bleu)}}
	\label{fig:reevaluation-base-regles-cs2-evaluation}

	\vspace{1cm}

	\begin{subfigure}[t]{0.1\textwidth}
		\centering
		\includegraphics[width=\textwidth]{cs2-1-reevaluation-square}
		\caption{}
		\label{fig:reevaluation-base-regles-cs2-1-reevaluation-square}
	\end{subfigure}
	\hfill
	\begin{subfigure}[t]{0.1\textwidth}
		\centering
		\includegraphics[width=\textwidth]{cs2-2-reevaluation-faceExtrusion}
		\caption{}
		\label{fig:reevaluation-base-regles-cs2-2-reevaluation-extrusion}
	\end{subfigure}
	\hfill
	\begin{subfigure}[t]{0.1\textwidth}
		\centering
		\includegraphics[width=\textwidth]{cs2-3-reevaluation-faceCollapse}
		\caption{}
		\label{fig:reevaluation-base-regles-cs2-3-reevaluation-collapse}
	\end{subfigure}
	\hfill
	\hspace{0.1\textwidth}
	\hfill
	\begin{subfigure}[t]{0.1\textwidth}
		\centering
		\includegraphics[width=\textwidth]{cs2-5-reevaluation-triangulation}
		\stepcounter{subfigure}\caption{}
		\label{fig:reevaluation-base-regles-cs2-5-reevaluation-triangulation}
	\end{subfigure}
	\hfill
	\hspace{0.1\textwidth}
	% \begin{subfigure}[t]{0.1\textwidth}
	% \centering
	% 	\includegraphics[width=\textwidth]{cs2-5-reevaluation-triangulation}
	% 	\caption{}
	% 	\label{fig:reevaluation-base-regles-cs2-CANCEL-6-reevaluation-color}
	% \end{subfigure}
	\hfill
	\hspace{0.1\textwidth}
	% \begin{subfigure}[t]{0.1\textwidth}
	% \centering
	% 	\includegraphics[width=\textwidth]{cs2-7-reevaluation-colourFaceBlue}
	% 	\caption{}
	% 	\label{fig:reevaluation-base-regles-cs2-7-reevaluation-color}
	% \end{subfigure}
	\caption{Réévaluation cas nº~2~:
		(\subref{fig:reevaluation-base-regles-cs2-1-reevaluation-square})
		\texttt{1-carre(pos)}~;
		(\subref{fig:reevaluation-base-regles-cs2-2-reevaluation-extrusion})
		\texttt{2-extrusion(PN\textsubscript{6},vec)}~;
		(\subref{fig:reevaluation-base-regles-cs2-3-reevaluation-collapse})
		\texttt{3-\allowbreak{}contraction\-Face(PN\textsubscript{47})}~;
		\texttt{DELETE~4-chanfrein(PN\textsubscript{40})}~;
		(\subref{fig:reevaluation-base-regles-cs2-5-reevaluation-triangulation})
		\texttt{5-\allowbreak{}triangulation\-Face(PN\textsubscript{33})}~;
		\texttt{6-coloration(PN\textsubscript{12},Rouge)}~;
		\texttt{7-coloration(PN\textsubscript{57},Bleu)}}
	\label{fig:reevaluation-base-regles-cs2-reevaluation}
\end{sidewaysfigure}

La seconde construction est illustrée
figure~\ref{fig:reevaluation-base-regles-cs2-evaluation}. %
Celle-ci va nous permettre de préciser la définition de \textit{stratégies
	de réévaluation}, notamment pour la suppression d'opérations. %
Nous détaillons, ci-après, sa spécification paramétrique contenue dans le
listing~\ref{lst:reevaluation-regles-specification-parametrique-init-cas-2}. %

% \xavier{Attention :\\
% 1) même remarque concernant le nom des fonctions\\
% 2) cette fois-ci, tu précises aussi le \textbf{numéro} des fonctions, par exemple "1-carre" au lieu de "carre" précédemment $\rightarrow$ sois homogène et précise \textbf{partout} le numéro, ou bien nulle part (et si c'est partout, simplifie ton texte : ce n'est pas la peine de préciser deux fois le numéro d'une opération avec le texte du style "La \textbf{troisième} application \textbf{3}-contraction[...])."
% }

La construction de cet objet est semblable à la construction nº~1. %
Le processus de modélisation commence avec l'application \texttt{1-\allowbreak{}creation\-Carre}
qui crée une face. %
L'application suivante \texttt{2-\allowbreak{}extrusion\-Face} extrude la face en un cube. %
La troisième application \texttt{3-\allowbreak{}contraction\-Face} contracte la face supérieure du
cube en un sommet. %
Ce sommet est ensuite chanfreiné par la quatrième application
\texttt{4-\allowbreak{}chanfrein\-Sommet}. %
L'application \texttt{5-\allowbreak{}triangulation\-Face} triangule la face frontale de l'objet. %
Enfin, les applications \texttt{6-\allowbreak{}coloration\-Face} et \texttt{7-\allowbreak{}coloration\-Face} colorient respectivement
deux faces de l'objet en rouge et en bleu. %
Notons que la face coloriée en rouge est la face issue du chanfreinage
(\texttt{4-\allowbreak{}chanfrein\-Sommet}). %

\noindent
\begin{minipage}{\textwidth} %% Je protège le listing qui apparaît en fin de page, sinon il est coupé par la figure suivante
\begin{lstlisting}[escapechar=\%,
style=script-style,
caption=Spécification paramétrique initiale du cas nº~2,
label=lst:reevaluation-regles-specification-parametrique-init-cas-2,
captionpos=b]
1-%\allowbreak{}creation\-Carre()%
2-%\allowbreak{}extrusion\-Face(\pndart{6}{1n6})%
3-%\allowbreak{}contraction\-Face(\pndart{47}{1n5; 2n6})%
4-%\allowbreak{}chanfrein\-Sommet(\pndart{40}{1n4; 2n4; 3n2})%
5-%\allowbreak{}triangulation\-Face(\pndart{33}{1n3; 2n2})%
6-%\allowbreak{}coloration\-Face(\pndart{12}{1n4; 2n4; 3n2; 4n2}, Rouge)%
7-%\allowbreak{}coloration\-Face(\pndart{57}{1n4; 2n4; 3n2; 4n1; 5n2}, Bleu)%
\end{lstlisting}
\end{minipage}
La version éditée de la spécification paramétrique consiste à supprimer la
quatrième opération, comme l'illustre le listing
~\ref{lst:reevaluation-regles-specification-parametrique-edit-cas-2}. %
L'application \texttt{4-\allowbreak{}chanfreinSommet} ne doit donc pas être réévaluée. %
Les applications \texttt{6-\allowbreak{}coloration\-Face} et
\texttt{7-\allowbreak{}coloration\-Face}, quant à elles, sont indirectement annulées, car les faces qu'elles doivent colorer ne sont pas recréées ainsi que l'atteste la
figure~\ref{fig:reevaluation-base-regles-cs2-reevaluation}.
%Dans la version éditée de la spécification paramétrique
%(liste~\ref{lst:reevaluation-regles-specification-parametrique-edit-cas-2}),
%figure~\ref{fig:reevaluation-base-regles-cs2-reevaluation}. %
%Les différentes étapes de réévaluation sont illustrées dans la
%figure~\ref{fig:reevaluation-base-regles-cs2-reevaluation}. %
%Les aspects remarquables de cette construction et de sa réévaluation sont le
%fait que la suppression de l'application \texttt{4-chanfrein} a un impact direct
%sur l'application \texttt{6-coloration} et indirect sur l'application
%\texttt{7-coloration}. %
En effet, l'application \texttt{6-coloration\-Face} référence la face créée par
\texttt{4-\allowbreak{}chanfrein\-Sommet} qui n'existe plus dans la réévaluation. %
De la même manière, la face initialement coloriée par \texttt{7-\allowbreak{}coloration\-Face}
n'est plus créée lors de la triangulation et ne peut donc pas être coloriée dans
la réévaluation. %

\begin{lstlisting}[escapechar=\%,
style=script-style,
caption=Spécification paramétrique éditée du cas nº~2,
label=lst:reevaluation-regles-specification-parametrique-edit-cas-2,
captionpos=b]
1-%\allowbreak{}creation\-Carre()%
2-%\allowbreak{}extrusion\-Face(\pndart{6}{1n6})%
3-%\allowbreak{}contraction\-Face(\pndart{47}{1n5; 2n6})%
DELETE 4-%\allowbreak{}chanfrein\-Sommet(\pndart{40}{1n4; 2n4; 3n2})%
5-%\allowbreak{}triangulation\-Face(\pndart{33}{1n3; 2n2})%
6-%\allowbreak{}coloration\-Face(\pndart{12}{1n4; 2n4; 3n2; 4n2}, Rouge)%
7-%\allowbreak{}coloration\-Face(\pndart{57}{1n4; 2n4; 3n2; 4n1; 5n2}, Bleu)%
\end{lstlisting}
